%% ============================================================================
\section{Related Work}
\label{sec:lit}

The theory of runs begins with Diamond and Dybvig~\cite{diamonddybvig1983}, whose
demand-deposit model leaves the run \emph{probability} undetermined, and global-games
methods select a unique equilibrium~\cite{carlssonvandamme1993,morrisshin1998,goldsteinpauzner2005}.
From that literature we take only the selection result, in which the marginal agent holds a
diffuse belief over the aggregate withdrawal mass. We do not reproduce the uniqueness
argument, and our result does not depend on it.

A literature since 2020 applies run theory to stablecoins, building on design
taxonomies~\cite{Moin_2020,Klages_Mundt_2020} and peg-arbitrage
mechanisms~\cite{Lyons_2023}. Ma, Zeng, and Zhang~\cite{mzz2025} build a global-games
stablecoin model, finding USDC's arbitrage less centralized than USDT's. Anadu et
al.~\cite{anadu2024} analyze the March 2023 de-peg via a discrete ``break-the-buck'' threshold.
Watsky et al.~\cite{watsky2024} give the institutional reason for that threshold, a closed
redemption window. USDC's primary redemption window
was open only to Circle's approved customers and was constrained by U.S. banking hours over
the March 10--13 weekend. Secondary prices therefore de-pegged while the arbitrage loop
was closed. A Federal Reserve post-mortem co-authored by Watsky reports the same
mechanism~\cite{fedshadow2025}. On AMM liquidity around the SVB
failure, Hernandez Cruz et al.~\cite{hernandezcruz2024} run a difference-in-differences
comparing USDC with USDT. Uhlig~\cite{uhlig2022} finds many TerraUSD holders behaved as
rational Bayesian updaters in the May 2022 TerraUSD collapse, which is evidence that
rational and panic behavior are hard to separate without a structural restriction. Liu,
Makarov, and Schoar~\cite{liumakarovschoar2023} reconstruct that collapse wallet-by-wallet.
Contemporaneous work surveys stablecoins as a payment technology and the financial risks
they pose~\cite{ahmed2025stablecoins}.

Other work measures the March~2023 de-peg without an expected-loss floor. Galati and
Capalbo~\cite{galaticapalbo2023} document the SVB link earliest, Diop, Chevallier, and
Sanhaji~\cite{diop2024fintech} rank USDC as most exposed in a daily cross-coin panel,
Cintra and Holloway~\cite{cintraholloway2023depegs} flag the de-peg about five hours before USDC crosses
\$0.99, Li et al.~\cite{li2026settlementfriction} frame it as settlement friction, and
Sankaewtong et al.~\cite{sankaewtong2026contagion} reconstruct the on-chain contagion
pathway at the transaction level. Our bound derives an expected loss from disclosed
exposure, and none of these papers does.

\paragraph{Contribution boundary.} Our work complements both lines of this literature. Structural
calibrations~\cite{mzz2025} estimate a run \emph{probability}, and event
studies~\cite{anadu2024} document \emph{flows}. We instead bound the \emph{realized
non-fundamental share} of one episode from an accounting identity. Two papers appear to conflict with this reading. Anadu et al.'s~\cite{anadu2024} \emph{rational} run concerns the
redemption decision. It does not concern whether the price was justified. Watsky et
al.'s~\cite{watsky2024} weekend discount is the microstructure friction the non-fundamental
share includes. Ahmed, Aldasoro, and Duley~\cite{ahmedaldasoroduley2024} run
synthetic-control tests on four \emph{other} episodes. They treat March 2023 only as a descriptive price path. Two further papers
address adjacent questions. Zhu~\cite{zhu2026redemptions} retains both large-sale and
collateral risk, splitting the space of fundamental states into regions by dominant risk
type, which classifies \emph{which} channel moves price and leaves \emph{how far}
unmeasured. Molinet and Filos-Ratsikas~\cite{molinetfilos2026} define fiat-backed coins as
the degenerate always-stable case ($1/p_t=1$), so they never analyze March 2023.

